% Bản thảo tiếng Việt — bản dịch làm việc của paper/manuscript_template.tex (bản tiếng Anh là bản nộp).
% Sinh bởi paper/build.py từ các run đã đăng ký; chỉ sửa tệp mẫu này hoặc sections/, không sửa manuscript.tex.
\documentclass[a4paper,preprint,12pt,authoryear]{elsarticle}
\usepackage{amsmath,amssymb}
\usepackage{booktabs}
\usepackage{graphicx}
\usepackage[table]{xcolor}
\usepackage[hidelinks,hypertexnames=false]{hyperref}
\usepackage[capitalise,noabbrev]{cleveref}
\usepackage{microtype}
\usepackage{float}
\setlength{\emergencystretch}{3em}
\pdfpagewidth=\paperwidth \pdfpageheight=\paperheight
\crefname{section}{Mục}{Mục}\Crefname{section}{Mục}{Mục}
\crefname{subsection}{Mục}{Mục}\Crefname{subsection}{Mục}{Mục}
\crefname{figure}{Hình}{Hình}\Crefname{figure}{Hình}{Hình}
\crefname{table}{Bảng}{Bảng}\Crefname{table}{Bảng}{Bảng}
\crefname{appendix}{Phụ lục}{Phụ lục}\Crefname{appendix}{Phụ lục}{Phụ lục}
\renewcommand{\appendixname}{Phụ lục}
\abstracttitle{Tóm tắt}
\keywordtitle{Từ khóa}
\AtBeginDocument{\renewcommand{\figurename}{Hình}\renewcommand{\tablename}{Bảng}\renewcommand{\refname}{Tài liệu tham khảo}}
\makeatletter
\g@addto@macro\appendix{\@addtoreset{table}{section}\@addtoreset{figure}{section}}
\makeatother
\myfooter[L]{Bản dịch tiếng Việt — bản nộp là bản tiếng Anh}
\journal{Comput. Educ.: Artif. Intell.}

\begin{document}
\begin{frontmatter}

\title{Khi test trượt là chưa đủ: đánh giá phân cụm quan niệm sai lầm trong bài lập trình C bằng nhãn từ bản sửa của chính sinh viên}

\begin{abstract}
%%INCLUDE sections/abstract.tex
\end{abstract}

\begin{keyword}
giáo dục lập trình \sep quan niệm sai lầm \sep phân cụm \sep sửa lỗi chương trình \sep quy nạp luật \sep khai phá dữ liệu giáo dục
\end{keyword}

\end{frontmatter}

%%INCLUDE sections/intro.tex
%%INCLUDE sections/related.tex
%%INCLUDE sections/methods.tex
%%INCLUDE sections/results.tex
%%INCLUDE sections/discussion.tex
%%INCLUDE sections/conclusion.tex

\appendix
%%INCLUDE sections/appendix.tex

\bibliographystyle{elsarticle-harv}
\bibliography{references}

\end{document}
